\section{Completing the arithmetic degree boundary}
\label{sec:composite-degree-completion}

The relative quantum registers admit a concrete completion across extension
degrees. Their finite comparison maps come from arithmetic field embeddings,
and their unnormalized traces come from the first nonzero copied-reference
coefficients. These two pieces of data determine the completion studied
here. Its Frobenius action remains a legitimate quantum automorphism and
becomes outer in the chosen minimal unitization. A separate positive degree
regularization supplies normalized states and an explicit zeta quotient.
The two statements in this section passed independent capped review;
the tail cutoff is explicitly restricted to integers.

\subsection{Divisor corners and their arithmetic provenance}

For an extension of relative degree $N$, the moving common observable
algebra consists of one matrix block $M_d(\mathbb C)$ for each divisor
$d>1$ of $N$. Its relative Frobenius is the cyclic shift
$S_d|k\rangle=|k+1\pmod d\rangle$ in each block. The randomized copied
reference has first coefficient $s\varphi(d)/d$ on that block, where
$s$ is the absolute degree of the base field. Dividing by this common
positive factor leaves the weights used below.
\provenance{D1602, D1603, D1607, CMP-REL, CMP-NATURAL, CMP-TRACE}
{relational-construction.md; boundary-and-spectrum.md}
{composite boundary checker}{composite-boundary-adjudication.md}

\begin{definition}[Divisor corners and a concrete degree completion]
For $N\geq2$ put
\[
 A_N=\bigoplus_{\substack{d\mid N\\d>1}}M_d(\mathbb C),\qquad
 \eta_d=\frac{\varphi(d)}d,\qquad
 \theta_N(a)=\sum_{\substack{d\mid N\\d>1}}\eta_d\operatorname{tr}_d(a_d),
 \quad \operatorname{tr}_d=\frac1d\operatorname{Tr}.
\]
Use the maximum of the block operator norms. For $N\mid M$, let
$j_{N,M}:A_N\to A_M$ extend block coordinates by zero.
On the Hilbert space
\[
 \mathcal H_{\rm deg}=\bigoplus_{d\geq2}^{\ell^2}\mathbb C^d,
\]
let $A_{\rm deg}$ be the block-diagonal algebra of sequences $(a_d)$ with
$\|a_d\|\to0$, in the supremum norm, and let
$B_{\rm deg}=A_{\rm deg}+\mathbb C I$ be its concrete minimal unitization.
For $a\geq0$ in $A_{\rm deg}$ define the extended-valued functional
\[
 \theta(a)=\sum_{d\geq2}\eta_d\operatorname{tr}_d(a_d).
\]
\end{definition}
\begin{scope}
The inclusions are nonunital corners. This is a specified completion of
the common arithmetic blocks; no total tensor or coherent direct-sum rule
on the distinguished primitive systems is stipulated. The Hilbert-space
symbol is not a Hamiltonian.
\end{scope}
\provenance{D1611}{definitions.md}{composite spectrum checker}
{composite-spectrum-adjudication.md}

The corner inclusion has an arithmetic interpretation. Fix a base
$K=\mathbb F_q$ and $N\mid M$. Every Frobenius orbit of length $d\mid N$
in $\mathbb F_{q^M}$ is already contained in $\mathbb F_{q^N}$: its labels
are fixed by the $N$th Frobenius power. Thus the relational encoder carries
the entire smaller period block to the entire same-period block upstairs.
Conjugation by that encoder, extending by zero off its image, realizes
$j_{N,M}$. In particular,
\[
 \theta_M\circ j_{N,M}=\theta_N.
\]
After normalizing the finite traces, restriction instead has success
$\theta_N(1)/\theta_M(1)$, with the smaller normalized reference as its
conditional state.

\subsection{Frobenius dynamics and positive degree states}

\begin{definition}[Completed Frobenius and degree regularization]
On $\mathcal H_{\rm deg}$ put
\[
 U_{\rm deg}=\bigoplus_{d\geq2}S_d,\qquad
 \alpha(a)_d=S_d a_d S_d^*,
\]
and let $\alpha$ fix the scalar unit of $B_{\rm deg}$.
For a real $\beta>1$, define
\begin{align*}
 Z_{\rm deg}(\beta)&=\sum_{d\geq2}\frac{\varphi(d)}{d^{\beta+1}},&
 D_\beta&=\bigoplus_{d\geq2}d^{-\beta}I_d,\\
 \rho_\beta&=\bigoplus_{d\geq2}
       \frac{\varphi(d)}{Z_{\rm deg}(\beta)d^{\beta+2}}I_d,&
 \omega_\beta(a)&=\operatorname{Tr}_{\mathcal H_{\rm deg}}(\rho_\beta a).
\end{align*}
Define $\zeta(x)=\sum_{m\geq1}m^{-x}$ for real $x>1$, and for an integer
$j$ define the implementing-operator moment
\[
 M_\beta(j)=\operatorname{Tr}_{\mathcal H_{\rm deg}}
                   (\rho_\beta U_{\rm deg}^{\,j}).
\]
\end{definition}
\begin{scope}
The regulator is additional central degree-weight data, not an arithmetic
amplitude or a Frobenius generator. The moment uses the named Hilbert
implementation; $U_{\rm deg}$ need not belong to $B_{\rm deg}$. No
analytic continuation or unregularized infinite trace is stipulated.
\end{scope}
\provenance{D1612}{definitions.md}{composite spectrum checker}
{composite-spectrum-adjudication.md}

\begin{theorem}[An arithmetic degree completion with outer Frobenius]
\statusProved\quad
The moving common algebras $A_N$ with the fixed-base arithmetic corner
maps form a directed system whose norm completion is $A_{\rm deg}$.
The compatible finite traces extend to the faithful, lower-semicontinuous
semifinite trace $\theta$ on $A_{\rm deg}$.

Relative Frobenius extends to the unital CP automorphism $\alpha$ of
$B_{\rm deg}$. This automorphism is outer in $B_{\rm deg}$ and is spatially
implemented by $U_{\rm deg}$ on $\mathcal H_{\rm deg}$. The implementing
unitary has all roots of unity as point spectrum, each with infinite
multiplicity, and the whole unit circle as its spectrum.
\end{theorem}
\begin{scope}
Outerness refers specifically to the minimal unitization. The same action
is inner in the larger product multiplier algebra. The trace is densely
defined on the nonunital algebra, not asserted densely defined on its
unitization. This is a discrete cycle dynamics, with no Riemann-zero
identification or rescaled continuous-time generator asserted.
\end{scope}
\provenance{CMP-COMPLETE}{degree-completion.md, Sections 1--2}
{composite spectrum checker}{composite-spectrum-adjudication.md}

\begin{proof}[Construction and the source of outerness]
Every finite collection of degrees lies among the divisors of a common
multiple. The algebraic union is therefore all finitely supported matrix
blocks, whose supremum-norm completion is precisely the sequence algebra
with vanishing tails. Conjugating each block by its arithmetic cyclic
shift preserves this norm and positivity at every matrix level.

If a unitary $w=\lambda I+a$ in the minimal unitization implemented
$\alpha$, then $\|a_d\|\to0$ and $|\lambda|=1$. For a rank-one block
projection $e_d=|0\rangle\langle0|$,
\[
 \|w_d e_d w_d^*-e_d\|\leq2\|w_d-\lambda I_d\|\longrightarrow0.
\]
But cyclic Frobenius sends $e_d$ to the orthogonal projection
$|1\rangle\langle1|$, whose difference from $e_d$ has norm one for
every $d\geq2$. This contradicts inner implementation in the unitization.

The positive trace is the supremum of its finite block sums. Finite
central cutoffs approximate positive elements in norm and in increasing
trace, proving the stated semifiniteness and lower semicontinuity.
Finally, finite cycle eigenvectors form an orthonormal basis of the Hilbert
sum. Roots of unity are dense on the circle; off it the block resolvents
have a uniform bound. This gives the stated spectrum.
\end{proof}

The regularized state is stationary, but the action is observable on
other preparations. In each inherited finite block, a relative-index
state at zero is carried to the orthogonal state at one. The coherent
three-register arithmetic test from the preceding section remains its
finite realization. No unrelated quantum block has been appended to
produce this action.

\begin{theorem}[A positive degree reference and its convergent zeta identity]
\statusProved\quad
For every real $\beta>1$, the operator $\rho_\beta$ is positive and trace
class with ordinary trace one. Its state on $B_{\rm deg}$ is faithful
and $\alpha$-invariant. For every integer cutoff $D\geq1$, the normalizer satisfies
\[
 Z_{\rm deg}(\beta)=\frac{\zeta(\beta)}{\zeta(\beta+1)}-1,
 \qquad
 \sum_{d>D}\frac{\varphi(d)}{d^{\beta+1}}
       \leq\frac{D^{1-\beta}}{\beta-1}\quad(D\geq1).
\]
For a nonzero integer $j$,
\[
 M_\beta(j)=\frac1{Z_{\rm deg}(\beta)}
       \sum_{\substack{d\mid |j|\\d\geq2}}
              \frac{\varphi(d)}{d^{\beta+1}},
 \qquad M_\beta(0)=1.
\]
\end{theorem}
\begin{scope}
The zeta identity is established in its convergent real domain. The
degree regulator is specified rather than claimed unique. The moment is
an implementing-operator trace paired with a trace-class state, not the
trace of a conjugation channel or an unregularized infinite Frobenius trace.
\end{scope}
\provenance{CMP-MELLIN}{degree-completion.md, Section 3}
{composite spectrum checker}{composite-spectrum-adjudication.md}

\begin{proof}[Positive coefficients and divisor convolution]
The bound $\varphi(d)\leq d$ compares the normalizer to the convergent
series $\sum d^{-\beta}$ and gives the integral tail estimate. Each
$d$-dimensional block contributes $\varphi(d)/(Z_{\rm deg}d^{\beta+1})$
to the ordinary trace of $\rho_\beta$, so these contributions sum to one.
Every eigenvalue is positive and each block is scalar, giving faithfulness
and Frobenius invariance.

Counting integers by their greatest common divisor with $k$ gives
$\sum_{d\mid k}\varphi(d)=k$. Absolute convergence then permits grouping
the product of positive series:
\[
 \left(\sum_{d\geq1}\frac{\varphi(d)}{d^{\beta+1}}\right)
 \zeta(\beta+1)
 =\sum_{k\geq1}\frac{\sum_{d\mid k}\varphi(d)}{k^{\beta+1}}
 =\zeta(\beta).
\]
Removing the $d=1$ term gives the claimed subtraction. The underlying
classical totient Dirichlet identity is also recorded in
\href{https://dlmf.nist.gov/27.4.E6}{NIST DLMF, equation 27.4.6}. The moment formula
uses $\operatorname{Tr}(S_d^j)=d$ when $d\mid j$, and zero otherwise.
The case $j=0$ uses trace-one normalization separately.
\end{proof}

The infinite unregularized trace does not follow by dropping the regulator:
the identity has infinite $\theta$-mass, so the unitary powers are not
$\theta$-integrable. Also, if logarithmic degree weights are written as
an operator with value $\log d$ on the $d$th block, that operator is central.
Its conjugation flow fixes the block algebra and must not be identified
with the nontrivial relative Frobenius action.

This completion gives one explicit positive quantum realization of the
surviving degree data. It preserves the arithmetic cyclic action while
leaving the fate of other composition laws open. For example, two
independent $2$-cycles have two $2$-cycles under simultaneous action;
they are not a single $4$-cycle. The corresponding matrix algebras can
have the same dimension while their distinguished Frobenius actions differ.
The existing greatest-common-divisor/least-common-multiple reblocking keeps
that multiplicity information when such independent composites are retained.
\provenance{FRL-COMP}{orbit-composition.md}{Frobenius boundary checker}
{frobenius-boundary-adjudication.md}
